%! Introduction to Eigenvalues — Unit 6, Lesson 6.1 — Exit Ticket (Answer Key)
\documentclass[10pt]{article}
\usepackage{linalg-article}
\usepackage{linalg-key}   % pulls in -boxes; do NOT also load -boxes
\newcommand{\TallMath}[1]{$\displaystyle #1\rule[-1.4em]{0pt}{3.2em}$}
\newcommand{\xx}{\mathbf{x}}
\newcommand{\zero}{\mathbf{0}}

\begin{document}

\pageheader{Unit 6, Lesson 6.1}{Exit Ticket --- Answer Key}
\namedateperiod

\vspace{0.15in}

No notes. Show your reasoning. Use $A=\begin{bmatrix} 2 & 2 \\ 1 & 3 \end{bmatrix}$ for items 1--2.

\begin{enumerate}[leftmargin=1.8em, itemsep=14pt]
  \item \textbf{Find the eigenvalues.} Form $A-\lambda I$, compute $\det(A-\lambda I)$, set it to $0$,
        and solve for both eigenvalues $\lambda$.
        \ansline{$A-\lambda I=\begin{bsmallmatrix}2-\lambda&2\\1&3-\lambda\end{bsmallmatrix}$; $\det=(2-\lambda)(3-\lambda)-2=\lambda^2-5\lambda+4=(\lambda-1)(\lambda-4)$, so $\lambda=1$ or $\lambda=4$.}
  \item \textbf{Find an eigenvector.} For the \emph{larger} eigenvalue, solve $(A-\lambda I)\xx=\zero$
        to find one eigenvector.
        \ansline{$\lambda=4$: $A-4I=\begin{bsmallmatrix}-2&2\\1&-1\end{bsmallmatrix}\Rightarrow x_1=x_2\Rightarrow\xx=\begin{bsmallmatrix}1\\1\end{bsmallmatrix}$. (Check: $A\begin{bsmallmatrix}1\\1\end{bsmallmatrix}=\begin{bsmallmatrix}4\\4\end{bsmallmatrix}=4\begin{bsmallmatrix}1\\1\end{bsmallmatrix}$.)}
  \item \textbf{Explain.} Why does setting $\det(A-\lambda I)=0$ find the eigenvalues? (One sentence ---
        think about what $(A-\lambda I)\xx=\zero$ needs.)
        \ansline{An eigenvector is a \emph{nonzero} $\xx$ with $(A-\lambda I)\xx=\zero$; a matrix crushes a nonzero vector to $\zero$ only when it is singular, and singular means $\det(A-\lambda I)=0$.}
\end{enumerate}

\begin{teachernote}
Sort into ``got it / arithmetic slip / concept slip.'' Item~1 is Step~1 ($\lambda=1,4$); item~2 is
Step~2 (eigenvector $\begin{bsmallmatrix}1\\1\end{bsmallmatrix}$ for $\lambda=4$); item~3 is the
\emph{why} (nonzero nullspace $\Rightarrow$ singular $\Rightarrow\det=0$). Watch for a sign slip in
$(2-\lambda)(3-\lambda)$ and for students who find the eigenvalues but skip solving
$(A-\lambda I)\xx=\zero$. If many miss item~3, re-anchor on the Unit~5 fact ``singular
$\Leftrightarrow\det=0$'' before Lesson~6.2.
\end{teachernote}

\end{document}
